\documentclass{article}
\usepackage[UTF8]{ctex}
\usepackage{amsmath,amssymb,amsfonts}
\usepackage{enumitem}
\usepackage{geometry}
\geometry{a4paper, margin=1in}
% === 错题本排版优化（不参与题目标识） ===
\setlength{\parindent}{0pt}
\setlength{\parskip}{0.35em plus 0.10em minus 0.08em}
\setlength{\emergencystretch}{3em}
\setlist{itemsep=0.25em, topsep=0.40em, parsep=0pt}
\allowdisplaybreaks[2]
% === 错题本排版优化结束 ===
\begin{document}
\section*{广西大学2026年数学分析真题}
\section*{164．广西大学2026年数学分析试卷}

\subsection{题目1}
1．已知点集 \(S \subset \mathbb{R}\) ，请叙述聚点的定义，若 \(S\) 为有限点集，它是否存在聚点，并给出理由。

\subsection{题目2}
2．已知 \(I\) 为函数 \(f(x)\) 定义域中的一个开区间，请用 \(\varepsilon-\delta\) 语言叙述 \(f\) 在 \(I\) 内不连续的定义，并给出所有间断点的类型及定义。

\subsection{题目3}
3．请叙述二元函数 \(f(x, y)\) 在点 \(P_{0}\left(x_{0}, y_{0}\right)\) 处可微的定义及在该点的全微分，令 \(z_{0}=f\left(x_{0}, y_{0}\right)\) ，请写出 \(f(x, y)\) 在点 \(P_{0}\left(x_{0}, y_{0}\right)\) 处可微时，曲面 \(z=f(x, y)\) 在点 \(P\left(x_{0}, y_{0}, z_{0}\right)\) 处的切平面方程．

\subsection{题目4}
1．设 \(\left\{\begin{array}{l}x(t)=\int_{0}^{t}\left(e^{u}-1\right) \mathrm{d} u, \\ y(t)=\ln \left(1+t^{2}\right) .\end{array}\right.\) 求 \(\lim _{t \rightarrow 0} \frac{y \sin y}{x^{2}}\) ．

\subsection{题目5}
2．求定积分 \(I_{n}=\int_{0}^{\frac{\pi}{2}} \cos ^{n} x \mathrm{~d} x(n=1,2, \cdots)\) ．

\subsection{题目6}
1．设 \(a_{1}, a_{2}, \cdots, a_{m}\) 为 \(m\) 个正数，证明：\(\sqrt[n]{a_{1}^{n}+a_{2}^{n}+\cdots+a_{m}^{n}}\) 极限存在，并求极限值．

\subsection{题目7}
2．请叙述有限覆盖定理并用其证明：若函数 \(f(x)\) 在 \([a, b]\) 上连续，则 \(f(x)\) 在 \([a, b]\) 上一致连续．

\subsection{题目8}
3．请叙述罗尔定理和拉格朗日中值定理，已知函数 \(f(x)\) 在区间 \([a, b]\) 上满足拉格朗日中值定理的条件，\(S(x)\) 表示以 \((a, f(a)),(b, f(b)),(x, f(x))\) 三点为顶点构成的三角形的面积，请利用 \(S(x)\) 和罗尔定理证明拉格朗日中值定理．

\subsection{题目9}
4．讨论函数 \(\Gamma(p)=\int_{0}^{+\infty} x^{p-1} e^{-x} \mathrm{~d} x\) 的收敛性并证明．

\subsection{题目10}
5．设函数 \(f(x, y) \in C^{2}\left(\mathbb{R}^{2}\right)\) ，对任意的 \((x, y) \in \mathbb{R}^{2}\) ，都有 \(f_{x x}(x, y)+f_{y y}(x, y)<0\) ，请写出 \(f(x, y)\) 的黑塞（Hesse）矩阵，并证明：\(f(x, y)\) 在 \(\mathbb{R}^{2}\) 上没有极小值点．

\subsection{题目11}
6．证明：级数 \(\frac{1!}{1}+\frac{2!}{2^{2}}+\frac{3!}{3^{3}}+\cdots+\frac{n!}{n^{n}}+\cdots\) 收敛。
\end{document}
